\chapter{The volume spectrum and Weyl law}\label{sec:volume-spectrum}
In this section $M^{n+1}$ is smooth, closed, and connected, and $n\geq1$. We use mod-two boundary cycles with the flat topology. The topology of this component provides a distinguished class $\bar\lambda\in H^1(\mathcal Z_n^0(M;\mathbb Z_2);\mathbb Z_2)$: the nonzero generator whose cup powers detect families. Its identification via Almgren's cycle-space theory is an external topological input; see \cite{MNIndex}, Section 5. The supplied published Weyl-law paper \cite{LMN} uses flat-continuous families with no concentration of mass.

\begin{definition}[No concentration of mass]\label{def:no-concentration}
A family $\Phi:X\to\mathcal Z_n^0(M;\mathbb Z_2)$ has no concentration of mass if
\[
 \lim_{r\downarrow0}\sup_{x\in X,\,p\in M}
 \|\Phi(x)\|(B_r(p))=0.
\]
The supremum over all parameters is essential: the absence of atoms in each individual slice does not give this uniform condition. See \cite{LMN}, Section 2.6.
\end{definition}

\begin{definition}[Sweepouts and the volume spectrum]\label{def:p-width}
For a positive integer $p$, a $p$-sweepout is a flat-continuous map $\Phi:X\to\mathcal Z_n^0(M;\mathbb Z_2)$ from a finite cubical complex such that
$\Phi^*(\bar\lambda^p)\neq0$ in $H^p(X;\mathbb Z_2)$.
In the present convention admissible sweepouts also have no concentration of mass and bounded maximal mass. Define
\[
 \omega_p(M,g)=\inf_\Phi\sup_{x\in X}\mathbf M_g(\Phi(x)).
\]
The infimum permits different parameter complexes and different homotopy classes. Thus $\omega_p$ is not the width of an arbitrary fixed class $\Pi$. We use the boundary component explicitly, as in \cite{MNIndex}, Sections 1 and 5, with the admissibility condition of \cite{LMN}, Section 2.6. Comparison with the $\mathbf F$-continuous convention uses interpolation, an external dependency.
\end{definition}

\begin{lemma}[Monotonicity in the sweepout index]\label{lem:width-monotonicity}
For every $p\geq1$, $\omega_p(M,g)\leq\omega_{p+1}(M,g)$.
\end{lemma}
\begin{proof}
If $\Phi^*(\bar\lambda^{p+1})\neq0$, then $\Phi^*(\bar\lambda^p)$ cannot vanish: multiplying a zero class by $\Phi^*(\bar\lambda)$ would give zero. Hence every admissible $(p+1)$-sweepout is an admissible $p$-sweepout. Taking the infima in Definition \ref{def:p-width} gives the inequality. This argument gives no strict inequality.
\end{proof}

\begin{lemma}[Scaling of widths]\label{lem:width-scaling}
For $r>0$ and $g_r=r^2g$,
\[
 \omega_p(M,g_r)=r^n\omega_p(M,g),\qquad
 \operatorname{Vol}(M,g_r)=r^{n+1}\operatorname{Vol}(M,g).
\]
\end{lemma}
\begin{proof}
The identity on $M$ preserves the flat topology, the component of boundary cycles, and the detecting class. An $n$-dimensional mass is multiplied by $r^n$; balls of radius $s$ for $g_r$ are balls of radius $s/r$ for $g$. Thus the no-concentration condition is preserved, and the admissible families are the same. Take their maximal masses and then their infimum. The volume identity follows by scaling the $(n+1)$-dimensional volume element.
\end{proof}

\begin{theorem}[Weyl law for the volume spectrum]\label{thm:weyl-volume}
For each $n\geq1$ there is a constant $a(n)>0$, depending only on $n$, such that every smooth closed connected $(M^{n+1},g)$ satisfies
\[
 \lim_{p\to\infty}p^{-1/(n+1)}\omega_p(M,g)
 =a(n)\operatorname{Vol}(M,g)^{n/(n+1)}.
\]
This is the closed-manifold case of \cite{LMN}, Theorem 1.1, p. 934. The source also treats compact manifolds with boundary using relative cycles. The proof, based on localization and comparison with Euclidean domains, is not yet included. The exponent agrees with Lemma \ref{lem:width-scaling}; the theorem does not assert an exact formula for any fixed finite $p$.
\end{theorem}

\chapter{Multiplicity and the Morse index}\label{sec:minmax-index}
The coefficients of a sweepout do not determine the multiplicities of its varifold limit. Index estimates require their own hypotheses in addition to the existence and regularity theory.

\begin{definition}[Bumpy metrics and the index of a minimal cycle]\label{def:bumpy-minimal-cycle}
A smooth metric is bumpy if no smooth closed immersed minimal hypersurface admits a nonzero normal Jacobi field. For an embedded minimal cycle
$V=\sum_j m_j|\Sigma_j|$ with distinct disjoint components, define
$\operatorname{index}(V)=\sum_j\operatorname{index}(\Sigma_j)$;
the integers $m_j$ do not multiply the index. Multiplicity one means that every $m_j$ is one. This convention is used in \cite{MNIndex}, Section 1. A disconnected hypersurface is allowed.
\end{definition}

\begin{theorem}[Index equality under a critical-set multiplicity hypothesis]\label{thm:minmax-index-equality}
Let $g$ be a smooth bumpy metric on a closed connected $M^{n+1}$ with $2\leq n\leq6$. Use the $\mathbf F$-continuous sweepout convention of \cite{MNIndex}. Let $\Pi$ be a homotopy class of $p$-sweepouts on a fixed finite $p$-dimensional simplicial complex, with $\mathbf L(\Pi)=\omega_p(M,g)$, and let $\Phi_i$ be a minimizing sequence in $\Pi$. Suppose that every embedded minimal cycle in its critical set has multiplicity one. Then that critical set contains an embedded minimal cycle $\Sigma$ satisfying
\[
 \operatorname{Area}_g(\Sigma)=\omega_p(M,g),\qquad
 \operatorname{index}(\Sigma)=p.
\]
Source: \cite{MNIndex}, Main Theorem 1.3, p. 5. The proof remains external. The multiplicity assumption concerns the whole indicated critical set, not merely one selected element; this theorem does not by itself prove multiplicity one.
\end{theorem}
